% =====================================================================
% Paper 16 — Provisional KZ-cosmology programme note (NOT a current
%            canonical GAP-4 closure paper)
% Status: [DRAFT] [NEEDS_UPDATE] [EXTERNAL_USE_FORBIDDEN]
% AUDIT-FLAG: AUDIT-2026-05-02-Wave7-Aux-Epoch-Overclaim
%   (Math314-AddD extension covering Wave 1/4/5 papers Paper-09..16).
%   Hostile-referee audit (2026-05-02): identified THREE structural
%   problems requiring repair before any external use:
%   (1) PHENOMENOLOGY MISMATCH: paper claims "Omega_GW peak at f ~ 10^-16
%       Hz (pulsar-timing array band)", but PTA bands are nHz
%       (~10^-9 to 10^-7 Hz). The claim is internally inconsistent
%       (paper later cites SKA/IPTA sensitivity at 10^-9 Hz). The 10^-16
%       Hz prediction is below ALL existing GW observational bands
%       (LIGO/Virgo: kHz; LISA: mHz; PTA: nHz; CMB B-mode: Hubble-scale).
%   (2) CANONICAL F-GAP4 REDEFINED: paper redefines F-GAP4 as
%       "SKA/IPTA non-detection by 2026-12-31" — but the canonical
%       F-GAP4 shell per Math300/Math311 is a defect-mass verdict
%       window: lambda_min(Box_E) -> mu_defect ∈ [10^13, 10^17] GeV
%       with PASS/FAIL/DEFER on solver/grid certificates. This paper
%       silently changes the gate definition.
%   (3) "Paper 16 completes the GAP-4 closure" wording incompatible
%       with current canonical Stage-2 composite = T3 and GAP-4 = T3.
%   External use is FORBIDDEN at present canonical tier; recast as a
%   Provisional KZ-cosmology programme note.
% Compile: pdflatex Paper-16.tex (×2 for refs)
% Canonical sources: Math156, Math159, Math172, Math196, Math218-AddA,
%                    Math300, Math311 (canonical F-GAP4 shell),
%                    Math314-AddD
% =====================================================================
% --- TECT canonical preamble (see _shared/tect-preamble.tex) ----------
% =====================================================================
% TECT canonical LaTeX preamble (revtex4-2 / single-column / theorem-safe)
% =====================================================================
% Maintainer: Jusang Lee (jtkor@outlook.com)
% Binding from: 2026-05-06
% Purpose: a single source-of-truth preamble for every TECT paper
% (Paper-00 .. Paper-10 + Auxiliary notes). Designed to (a) eliminate
% font-corruption on pdflatex (T1 fontenc + lmodern), (b) make
% theorem-heavy mathematical-physics content render cleanly in a
% single-column layout (PRD class option, NOT PRL two-column), (c)
% provide consistent theorem numbering across all TECT papers via
% amsthm with a shared counter set, (d) make hyperref + cleveref
% cross-references resolve cleanly with the standard two-pass
% pdflatex compile.
%
% Usage in each Paper-NN.tex:
%   % =====================================================================
% TECT canonical LaTeX preamble (revtex4-2 / single-column / theorem-safe)
% =====================================================================
% Maintainer: Jusang Lee (jtkor@outlook.com)
% Binding from: 2026-05-06
% Purpose: a single source-of-truth preamble for every TECT paper
% (Paper-00 .. Paper-10 + Auxiliary notes). Designed to (a) eliminate
% font-corruption on pdflatex (T1 fontenc + lmodern), (b) make
% theorem-heavy mathematical-physics content render cleanly in a
% single-column layout (PRD class option, NOT PRL two-column), (c)
% provide consistent theorem numbering across all TECT papers via
% amsthm with a shared counter set, (d) make hyperref + cleveref
% cross-references resolve cleanly with the standard two-pass
% pdflatex compile.
%
% Usage in each Paper-NN.tex:
%   % =====================================================================
% TECT canonical LaTeX preamble (revtex4-2 / single-column / theorem-safe)
% =====================================================================
% Maintainer: Jusang Lee (jtkor@outlook.com)
% Binding from: 2026-05-06
% Purpose: a single source-of-truth preamble for every TECT paper
% (Paper-00 .. Paper-10 + Auxiliary notes). Designed to (a) eliminate
% font-corruption on pdflatex (T1 fontenc + lmodern), (b) make
% theorem-heavy mathematical-physics content render cleanly in a
% single-column layout (PRD class option, NOT PRL two-column), (c)
% provide consistent theorem numbering across all TECT papers via
% amsthm with a shared counter set, (d) make hyperref + cleveref
% cross-references resolve cleanly with the standard two-pass
% pdflatex compile.
%
% Usage in each Paper-NN.tex:
%   \input{../_shared/tect-preamble.tex}
%   \begin{document}
%   ...
%   \end{document}
%
% Build (every paper): `pdflatex Paper-NN && pdflatex Paper-NN`
% (the second pass resolves \ref{} / \cref{} forward references).
% A bash/PowerShell wrapper that automates this is at
% Docs/papers/papers/_shared/build-paper.sh / build-paper.ps1.
% =====================================================================

% ---- Document class --------------------------------------------------
% PRD class option of revtex4-2 with [reprint,onecolumn] gives the same
% APS heading / by-line / abstract style as PRL but in a wider single
% column that handles long display equations and theorem environments
% without column-break artefacts. nofootinbib keeps citations out of
% footnotes. The 11pt option restores standard font metrics; with T1
% fontenc + lmodern this gives non-bitmap glyphs.
\documentclass[%
  aps,prd,reprint,onecolumn,nofootinbib,11pt,%
  amsmath,amssymb,floatfix,longbibliography%
]{revtex4-2}

% ---- Encoding & fonts ------------------------------------------------
\usepackage[T1]{fontenc}        % proper 8-bit font encoding (no font corruption)
\usepackage[utf8]{inputenc}     % allow UTF-8 source
\usepackage{lmodern}            % scalable Latin Modern (replaces bitmap CM)
\usepackage{textcomp}           % extra text symbols
\usepackage{microtype}          % subtle kerning improvements

% ---- UTF-8 → LaTeX character mappings --------------------------------
% Maps the Unicode characters that occur in TECT body text (legacy from
% the chat-content auto-archival rule, CLAUDE.md §4) to LaTeX-safe
% commands. Without these declarations, T1+inputenc rejects the chars
% with "Unicode character X not set up for use with LaTeX".
% Adding declarations centrally (here) is preferable to per-file edits.
\DeclareUnicodeCharacter{00A7}{\S}              % §  (section sign)
\DeclareUnicodeCharacter{00B2}{\ensuremath{^{2}}}        % ²
\DeclareUnicodeCharacter{00B3}{\ensuremath{^{3}}}        % ³
\DeclareUnicodeCharacter{00B5}{\ensuremath{\mu}}         % µ (micro)
\DeclareUnicodeCharacter{00B7}{\ensuremath{\cdot}}       % · (middle dot)
\DeclareUnicodeCharacter{00D7}{\ensuremath{\times}}      % ×
\DeclareUnicodeCharacter{00E4}{\"a}             % ä
\DeclareUnicodeCharacter{00E9}{\'e}             % é
\DeclareUnicodeCharacter{00F6}{\"o}             % ö
\DeclareUnicodeCharacter{00FC}{\"u}             % ü
\DeclareUnicodeCharacter{010C}{\v{C}}           % Č
\DeclareUnicodeCharacter{010D}{\v{c}}           % č
\DeclareUnicodeCharacter{0394}{\ensuremath{\Delta}}      % Δ
\DeclareUnicodeCharacter{03A3}{\ensuremath{\Sigma}}      % Σ
\DeclareUnicodeCharacter{03A9}{\ensuremath{\Omega}}      % Ω
\DeclareUnicodeCharacter{03B1}{\ensuremath{\alpha}}      % α
\DeclareUnicodeCharacter{03B2}{\ensuremath{\beta}}       % β
\DeclareUnicodeCharacter{03B3}{\ensuremath{\gamma}}      % γ
\DeclareUnicodeCharacter{03B4}{\ensuremath{\delta}}      % δ
\DeclareUnicodeCharacter{03B5}{\ensuremath{\varepsilon}} % ε
\DeclareUnicodeCharacter{03BA}{\ensuremath{\kappa}}      % κ
\DeclareUnicodeCharacter{03BB}{\ensuremath{\lambda}}     % λ
\DeclareUnicodeCharacter{03BC}{\ensuremath{\mu}}         % μ
\DeclareUnicodeCharacter{03BD}{\ensuremath{\nu}}         % ν
\DeclareUnicodeCharacter{03C0}{\ensuremath{\pi}}         % π
\DeclareUnicodeCharacter{03C1}{\ensuremath{\rho}}        % ρ
\DeclareUnicodeCharacter{03C3}{\ensuremath{\sigma}}      % σ
\DeclareUnicodeCharacter{03C4}{\ensuremath{\tau}}        % τ
\DeclareUnicodeCharacter{03C6}{\ensuremath{\varphi}}     % φ
\DeclareUnicodeCharacter{03C7}{\ensuremath{\chi}}        % χ
\DeclareUnicodeCharacter{03C8}{\ensuremath{\psi}}        % ψ
\DeclareUnicodeCharacter{03C9}{\ensuremath{\omega}}      % ω
\DeclareUnicodeCharacter{2013}{--}              % – (en dash)
\DeclareUnicodeCharacter{2014}{---}             % — (em dash)
\DeclareUnicodeCharacter{2018}{`}               % ‘
\DeclareUnicodeCharacter{2019}{'}               % ’
\DeclareUnicodeCharacter{201C}{``}              % “
\DeclareUnicodeCharacter{201D}{''}              % ”
\DeclareUnicodeCharacter{2070}{\ensuremath{^{0}}}        % ⁰
\DeclareUnicodeCharacter{2071}{\ensuremath{^{i}}}        % ⁱ
\DeclareUnicodeCharacter{2122}{\textsuperscript{TM}}     % ™
\DeclareUnicodeCharacter{2126}{\ensuremath{\Omega}}      % Ω (ohm sign)
\DeclareUnicodeCharacter{210F}{\ensuremath{\hbar}}       % ℏ
\DeclareUnicodeCharacter{2115}{\ensuremath{\mathbb{N}}}  % ℕ
\DeclareUnicodeCharacter{2102}{\ensuremath{\mathbb{C}}}  % ℂ
\DeclareUnicodeCharacter{211D}{\ensuremath{\mathbb{R}}}  % ℝ
\DeclareUnicodeCharacter{2124}{\ensuremath{\mathbb{Z}}}  % ℤ
\DeclareUnicodeCharacter{211A}{\ensuremath{\mathbb{Q}}}  % ℚ
\DeclareUnicodeCharacter{2192}{\ensuremath{\to}}         % →
\DeclareUnicodeCharacter{21D2}{\ensuremath{\Rightarrow}} % ⇒
\DeclareUnicodeCharacter{2200}{\ensuremath{\forall}}     % ∀
\DeclareUnicodeCharacter{2203}{\ensuremath{\exists}}     % ∃
\DeclareUnicodeCharacter{2208}{\ensuremath{\in}}         % ∈
\DeclareUnicodeCharacter{2209}{\ensuremath{\notin}}      % ∉
\DeclareUnicodeCharacter{2227}{\ensuremath{\wedge}}      % ∧
\DeclareUnicodeCharacter{2228}{\ensuremath{\vee}}        % ∨
\DeclareUnicodeCharacter{2229}{\ensuremath{\cap}}        % ∩
\DeclareUnicodeCharacter{222A}{\ensuremath{\cup}}        % ∪
\DeclareUnicodeCharacter{2245}{\ensuremath{\cong}}       % ≅
\DeclareUnicodeCharacter{2248}{\ensuremath{\approx}}     % ≈
\DeclareUnicodeCharacter{2260}{\ensuremath{\neq}}        % ≠
\DeclareUnicodeCharacter{2264}{\ensuremath{\leq}}        % ≤
\DeclareUnicodeCharacter{2265}{\ensuremath{\geq}}        % ≥
\DeclareUnicodeCharacter{22C5}{\ensuremath{\cdot}}       % ⋅
\DeclareUnicodeCharacter{2295}{\ensuremath{\oplus}}      % ⊕
\DeclareUnicodeCharacter{2297}{\ensuremath{\otimes}}     % ⊗

% ---- Mathematics & symbols ------------------------------------------
\usepackage{amsmath,amssymb,bm,mathrsfs,mathtools}
\usepackage{amsthm}             % theorem environments (load BEFORE hyperref)

% ---- Theorem environments (shared counter set, TECT canonical) -------
% All theorems / lemmas / propositions / corollaries / remarks share a
% single counter so cross-references read consistently across a paper.
% Definitions get a separate counter (definitions are numbered in their
% own sequence by editorial convention).
\newtheorem{theorem}{Theorem}
\newtheorem{lemma}[theorem]{Lemma}
\newtheorem{proposition}[theorem]{Proposition}
\newtheorem{corollary}[theorem]{Corollary}

\theoremstyle{remark}
\newtheorem{remark}[theorem]{Remark}
\newtheorem{example}[theorem]{Example}

\theoremstyle{definition}
\newtheorem{definition}[theorem]{Definition}
\newtheorem{conjecture}[theorem]{Conjecture}
\newtheorem{hypothesis}[theorem]{Hypothesis}

% ---- Tables / floats / graphics --------------------------------------
\usepackage{booktabs}           % \toprule / \midrule / \bottomrule
\usepackage{array}
\usepackage{graphicx}
\usepackage{enumitem}           % \begin{itemize}[leftmargin=...] options
\usepackage{hyphenat}           % \hyp{} for breakable hyphens in \texttt
\usepackage{seqsplit}           % \seqsplit{} for long unbreakable strings
\sloppy                         % allow stretchier inter-word spacing to
                                % reduce overfull \hbox warnings on long URLs
                                % and long \texttt tokens (paragraph-level).
\emergencystretch=3em           % last-resort hbox stretch budget

% ---- Cross-references (hyperref last among the standard packages,
%      cleveref AFTER hyperref) ----------------------------------------
\usepackage[%
  colorlinks=true,%
  linkcolor=blue,%
  citecolor=blue,%
  urlcolor=blue,%
  breaklinks=true,%
  bookmarks=true,%
  bookmarksnumbered=true%
]{hyperref}
\usepackage[capitalise,nameinlink]{cleveref}

% ---- TECT-canonical author block macros -----------------------------
% (defined here so every paper has the same author / affiliation style)
\newcommand{\tectauthor}{%
  \author{Jusang Lee}%
  \email{jtkor@outlook.com}%
  \affiliation{Independent researcher, Republic of Korea}%
}

% ---- TECT-canonical archive footer ----------------------------------
\newcommand{\tectarchivefooter}{%
  \par\medskip
  \noindent\small\textit{Canonical archive}:
  \url{https://github.com/TECT-OpenPhysics/TECT}.\par%
}

% ---- TECT TODO / AUDIT-FLAG / HONEST-SCOPE inline marks --------------
% Used in drafts to highlight pending audit items without disturbing
% the body text style. Suppress via \tectsuppressmarks (define empty).
\providecommand{\tectauditflag}[1]{\textcolor{red}{\textbf{[AUDIT-FLAG: #1]}}}
\providecommand{\tecttodo}[1]{\textcolor{orange}{\textbf{[TODO: #1]}}}
\providecommand{\tecthonestscope}[1]{\textit{Honest scope: #1.}}

% =====================================================================
% End of TECT canonical preamble
% =====================================================================

%   \begin{document}
%   ...
%   \end{document}
%
% Build (every paper): `pdflatex Paper-NN && pdflatex Paper-NN`
% (the second pass resolves \ref{} / \cref{} forward references).
% A bash/PowerShell wrapper that automates this is at
% Docs/papers/papers/_shared/build-paper.sh / build-paper.ps1.
% =====================================================================

% ---- Document class --------------------------------------------------
% PRD class option of revtex4-2 with [reprint,onecolumn] gives the same
% APS heading / by-line / abstract style as PRL but in a wider single
% column that handles long display equations and theorem environments
% without column-break artefacts. nofootinbib keeps citations out of
% footnotes. The 11pt option restores standard font metrics; with T1
% fontenc + lmodern this gives non-bitmap glyphs.
\documentclass[%
  aps,prd,reprint,onecolumn,nofootinbib,11pt,%
  amsmath,amssymb,floatfix,longbibliography%
]{revtex4-2}

% ---- Encoding & fonts ------------------------------------------------
\usepackage[T1]{fontenc}        % proper 8-bit font encoding (no font corruption)
\usepackage[utf8]{inputenc}     % allow UTF-8 source
\usepackage{lmodern}            % scalable Latin Modern (replaces bitmap CM)
\usepackage{textcomp}           % extra text symbols
\usepackage{microtype}          % subtle kerning improvements

% ---- UTF-8 → LaTeX character mappings --------------------------------
% Maps the Unicode characters that occur in TECT body text (legacy from
% the chat-content auto-archival rule, CLAUDE.md §4) to LaTeX-safe
% commands. Without these declarations, T1+inputenc rejects the chars
% with "Unicode character X not set up for use with LaTeX".
% Adding declarations centrally (here) is preferable to per-file edits.
\DeclareUnicodeCharacter{00A7}{\S}              % §  (section sign)
\DeclareUnicodeCharacter{00B2}{\ensuremath{^{2}}}        % ²
\DeclareUnicodeCharacter{00B3}{\ensuremath{^{3}}}        % ³
\DeclareUnicodeCharacter{00B5}{\ensuremath{\mu}}         % µ (micro)
\DeclareUnicodeCharacter{00B7}{\ensuremath{\cdot}}       % · (middle dot)
\DeclareUnicodeCharacter{00D7}{\ensuremath{\times}}      % ×
\DeclareUnicodeCharacter{00E4}{\"a}             % ä
\DeclareUnicodeCharacter{00E9}{\'e}             % é
\DeclareUnicodeCharacter{00F6}{\"o}             % ö
\DeclareUnicodeCharacter{00FC}{\"u}             % ü
\DeclareUnicodeCharacter{010C}{\v{C}}           % Č
\DeclareUnicodeCharacter{010D}{\v{c}}           % č
\DeclareUnicodeCharacter{0394}{\ensuremath{\Delta}}      % Δ
\DeclareUnicodeCharacter{03A3}{\ensuremath{\Sigma}}      % Σ
\DeclareUnicodeCharacter{03A9}{\ensuremath{\Omega}}      % Ω
\DeclareUnicodeCharacter{03B1}{\ensuremath{\alpha}}      % α
\DeclareUnicodeCharacter{03B2}{\ensuremath{\beta}}       % β
\DeclareUnicodeCharacter{03B3}{\ensuremath{\gamma}}      % γ
\DeclareUnicodeCharacter{03B4}{\ensuremath{\delta}}      % δ
\DeclareUnicodeCharacter{03B5}{\ensuremath{\varepsilon}} % ε
\DeclareUnicodeCharacter{03BA}{\ensuremath{\kappa}}      % κ
\DeclareUnicodeCharacter{03BB}{\ensuremath{\lambda}}     % λ
\DeclareUnicodeCharacter{03BC}{\ensuremath{\mu}}         % μ
\DeclareUnicodeCharacter{03BD}{\ensuremath{\nu}}         % ν
\DeclareUnicodeCharacter{03C0}{\ensuremath{\pi}}         % π
\DeclareUnicodeCharacter{03C1}{\ensuremath{\rho}}        % ρ
\DeclareUnicodeCharacter{03C3}{\ensuremath{\sigma}}      % σ
\DeclareUnicodeCharacter{03C4}{\ensuremath{\tau}}        % τ
\DeclareUnicodeCharacter{03C6}{\ensuremath{\varphi}}     % φ
\DeclareUnicodeCharacter{03C7}{\ensuremath{\chi}}        % χ
\DeclareUnicodeCharacter{03C8}{\ensuremath{\psi}}        % ψ
\DeclareUnicodeCharacter{03C9}{\ensuremath{\omega}}      % ω
\DeclareUnicodeCharacter{2013}{--}              % – (en dash)
\DeclareUnicodeCharacter{2014}{---}             % — (em dash)
\DeclareUnicodeCharacter{2018}{`}               % ‘
\DeclareUnicodeCharacter{2019}{'}               % ’
\DeclareUnicodeCharacter{201C}{``}              % “
\DeclareUnicodeCharacter{201D}{''}              % ”
\DeclareUnicodeCharacter{2070}{\ensuremath{^{0}}}        % ⁰
\DeclareUnicodeCharacter{2071}{\ensuremath{^{i}}}        % ⁱ
\DeclareUnicodeCharacter{2122}{\textsuperscript{TM}}     % ™
\DeclareUnicodeCharacter{2126}{\ensuremath{\Omega}}      % Ω (ohm sign)
\DeclareUnicodeCharacter{210F}{\ensuremath{\hbar}}       % ℏ
\DeclareUnicodeCharacter{2115}{\ensuremath{\mathbb{N}}}  % ℕ
\DeclareUnicodeCharacter{2102}{\ensuremath{\mathbb{C}}}  % ℂ
\DeclareUnicodeCharacter{211D}{\ensuremath{\mathbb{R}}}  % ℝ
\DeclareUnicodeCharacter{2124}{\ensuremath{\mathbb{Z}}}  % ℤ
\DeclareUnicodeCharacter{211A}{\ensuremath{\mathbb{Q}}}  % ℚ
\DeclareUnicodeCharacter{2192}{\ensuremath{\to}}         % →
\DeclareUnicodeCharacter{21D2}{\ensuremath{\Rightarrow}} % ⇒
\DeclareUnicodeCharacter{2200}{\ensuremath{\forall}}     % ∀
\DeclareUnicodeCharacter{2203}{\ensuremath{\exists}}     % ∃
\DeclareUnicodeCharacter{2208}{\ensuremath{\in}}         % ∈
\DeclareUnicodeCharacter{2209}{\ensuremath{\notin}}      % ∉
\DeclareUnicodeCharacter{2227}{\ensuremath{\wedge}}      % ∧
\DeclareUnicodeCharacter{2228}{\ensuremath{\vee}}        % ∨
\DeclareUnicodeCharacter{2229}{\ensuremath{\cap}}        % ∩
\DeclareUnicodeCharacter{222A}{\ensuremath{\cup}}        % ∪
\DeclareUnicodeCharacter{2245}{\ensuremath{\cong}}       % ≅
\DeclareUnicodeCharacter{2248}{\ensuremath{\approx}}     % ≈
\DeclareUnicodeCharacter{2260}{\ensuremath{\neq}}        % ≠
\DeclareUnicodeCharacter{2264}{\ensuremath{\leq}}        % ≤
\DeclareUnicodeCharacter{2265}{\ensuremath{\geq}}        % ≥
\DeclareUnicodeCharacter{22C5}{\ensuremath{\cdot}}       % ⋅
\DeclareUnicodeCharacter{2295}{\ensuremath{\oplus}}      % ⊕
\DeclareUnicodeCharacter{2297}{\ensuremath{\otimes}}     % ⊗

% ---- Mathematics & symbols ------------------------------------------
\usepackage{amsmath,amssymb,bm,mathrsfs,mathtools}
\usepackage{amsthm}             % theorem environments (load BEFORE hyperref)

% ---- Theorem environments (shared counter set, TECT canonical) -------
% All theorems / lemmas / propositions / corollaries / remarks share a
% single counter so cross-references read consistently across a paper.
% Definitions get a separate counter (definitions are numbered in their
% own sequence by editorial convention).
\newtheorem{theorem}{Theorem}
\newtheorem{lemma}[theorem]{Lemma}
\newtheorem{proposition}[theorem]{Proposition}
\newtheorem{corollary}[theorem]{Corollary}

\theoremstyle{remark}
\newtheorem{remark}[theorem]{Remark}
\newtheorem{example}[theorem]{Example}

\theoremstyle{definition}
\newtheorem{definition}[theorem]{Definition}
\newtheorem{conjecture}[theorem]{Conjecture}
\newtheorem{hypothesis}[theorem]{Hypothesis}

% ---- Tables / floats / graphics --------------------------------------
\usepackage{booktabs}           % \toprule / \midrule / \bottomrule
\usepackage{array}
\usepackage{graphicx}
\usepackage{enumitem}           % \begin{itemize}[leftmargin=...] options
\usepackage{hyphenat}           % \hyp{} for breakable hyphens in \texttt
\usepackage{seqsplit}           % \seqsplit{} for long unbreakable strings
\sloppy                         % allow stretchier inter-word spacing to
                                % reduce overfull \hbox warnings on long URLs
                                % and long \texttt tokens (paragraph-level).
\emergencystretch=3em           % last-resort hbox stretch budget

% ---- Cross-references (hyperref last among the standard packages,
%      cleveref AFTER hyperref) ----------------------------------------
\usepackage[%
  colorlinks=true,%
  linkcolor=blue,%
  citecolor=blue,%
  urlcolor=blue,%
  breaklinks=true,%
  bookmarks=true,%
  bookmarksnumbered=true%
]{hyperref}
\usepackage[capitalise,nameinlink]{cleveref}

% ---- TECT-canonical author block macros -----------------------------
% (defined here so every paper has the same author / affiliation style)
\newcommand{\tectauthor}{%
  \author{Jusang Lee}%
  \email{jtkor@outlook.com}%
  \affiliation{Independent researcher, Republic of Korea}%
}

% ---- TECT-canonical archive footer ----------------------------------
\newcommand{\tectarchivefooter}{%
  \par\medskip
  \noindent\small\textit{Canonical archive}:
  \url{https://github.com/TECT-OpenPhysics/TECT}.\par%
}

% ---- TECT TODO / AUDIT-FLAG / HONEST-SCOPE inline marks --------------
% Used in drafts to highlight pending audit items without disturbing
% the body text style. Suppress via \tectsuppressmarks (define empty).
\providecommand{\tectauditflag}[1]{\textcolor{red}{\textbf{[AUDIT-FLAG: #1]}}}
\providecommand{\tecttodo}[1]{\textcolor{orange}{\textbf{[TODO: #1]}}}
\providecommand{\tecthonestscope}[1]{\textit{Honest scope: #1.}}

% =====================================================================
% End of TECT canonical preamble
% =====================================================================

%   \begin{document}
%   ...
%   \end{document}
%
% Build (every paper): `pdflatex Paper-NN && pdflatex Paper-NN`
% (the second pass resolves \ref{} / \cref{} forward references).
% A bash/PowerShell wrapper that automates this is at
% Docs/papers/papers/_shared/build-paper.sh / build-paper.ps1.
% =====================================================================

% ---- Document class --------------------------------------------------
% PRD class option of revtex4-2 with [reprint,onecolumn] gives the same
% APS heading / by-line / abstract style as PRL but in a wider single
% column that handles long display equations and theorem environments
% without column-break artefacts. nofootinbib keeps citations out of
% footnotes. The 11pt option restores standard font metrics; with T1
% fontenc + lmodern this gives non-bitmap glyphs.
\documentclass[%
  aps,prd,reprint,onecolumn,nofootinbib,11pt,%
  amsmath,amssymb,floatfix,longbibliography%
]{revtex4-2}

% ---- Encoding & fonts ------------------------------------------------
\usepackage[T1]{fontenc}        % proper 8-bit font encoding (no font corruption)
\usepackage[utf8]{inputenc}     % allow UTF-8 source
\usepackage{lmodern}            % scalable Latin Modern (replaces bitmap CM)
\usepackage{textcomp}           % extra text symbols
\usepackage{microtype}          % subtle kerning improvements

% ---- UTF-8 → LaTeX character mappings --------------------------------
% Maps the Unicode characters that occur in TECT body text (legacy from
% the chat-content auto-archival rule, CLAUDE.md §4) to LaTeX-safe
% commands. Without these declarations, T1+inputenc rejects the chars
% with "Unicode character X not set up for use with LaTeX".
% Adding declarations centrally (here) is preferable to per-file edits.
\DeclareUnicodeCharacter{00A7}{\S}              % §  (section sign)
\DeclareUnicodeCharacter{00B2}{\ensuremath{^{2}}}        % ²
\DeclareUnicodeCharacter{00B3}{\ensuremath{^{3}}}        % ³
\DeclareUnicodeCharacter{00B5}{\ensuremath{\mu}}         % µ (micro)
\DeclareUnicodeCharacter{00B7}{\ensuremath{\cdot}}       % · (middle dot)
\DeclareUnicodeCharacter{00D7}{\ensuremath{\times}}      % ×
\DeclareUnicodeCharacter{00E4}{\"a}             % ä
\DeclareUnicodeCharacter{00E9}{\'e}             % é
\DeclareUnicodeCharacter{00F6}{\"o}             % ö
\DeclareUnicodeCharacter{00FC}{\"u}             % ü
\DeclareUnicodeCharacter{010C}{\v{C}}           % Č
\DeclareUnicodeCharacter{010D}{\v{c}}           % č
\DeclareUnicodeCharacter{0394}{\ensuremath{\Delta}}      % Δ
\DeclareUnicodeCharacter{03A3}{\ensuremath{\Sigma}}      % Σ
\DeclareUnicodeCharacter{03A9}{\ensuremath{\Omega}}      % Ω
\DeclareUnicodeCharacter{03B1}{\ensuremath{\alpha}}      % α
\DeclareUnicodeCharacter{03B2}{\ensuremath{\beta}}       % β
\DeclareUnicodeCharacter{03B3}{\ensuremath{\gamma}}      % γ
\DeclareUnicodeCharacter{03B4}{\ensuremath{\delta}}      % δ
\DeclareUnicodeCharacter{03B5}{\ensuremath{\varepsilon}} % ε
\DeclareUnicodeCharacter{03BA}{\ensuremath{\kappa}}      % κ
\DeclareUnicodeCharacter{03BB}{\ensuremath{\lambda}}     % λ
\DeclareUnicodeCharacter{03BC}{\ensuremath{\mu}}         % μ
\DeclareUnicodeCharacter{03BD}{\ensuremath{\nu}}         % ν
\DeclareUnicodeCharacter{03C0}{\ensuremath{\pi}}         % π
\DeclareUnicodeCharacter{03C1}{\ensuremath{\rho}}        % ρ
\DeclareUnicodeCharacter{03C3}{\ensuremath{\sigma}}      % σ
\DeclareUnicodeCharacter{03C4}{\ensuremath{\tau}}        % τ
\DeclareUnicodeCharacter{03C6}{\ensuremath{\varphi}}     % φ
\DeclareUnicodeCharacter{03C7}{\ensuremath{\chi}}        % χ
\DeclareUnicodeCharacter{03C8}{\ensuremath{\psi}}        % ψ
\DeclareUnicodeCharacter{03C9}{\ensuremath{\omega}}      % ω
\DeclareUnicodeCharacter{2013}{--}              % – (en dash)
\DeclareUnicodeCharacter{2014}{---}             % — (em dash)
\DeclareUnicodeCharacter{2018}{`}               % ‘
\DeclareUnicodeCharacter{2019}{'}               % ’
\DeclareUnicodeCharacter{201C}{``}              % “
\DeclareUnicodeCharacter{201D}{''}              % ”
\DeclareUnicodeCharacter{2070}{\ensuremath{^{0}}}        % ⁰
\DeclareUnicodeCharacter{2071}{\ensuremath{^{i}}}        % ⁱ
\DeclareUnicodeCharacter{2122}{\textsuperscript{TM}}     % ™
\DeclareUnicodeCharacter{2126}{\ensuremath{\Omega}}      % Ω (ohm sign)
\DeclareUnicodeCharacter{210F}{\ensuremath{\hbar}}       % ℏ
\DeclareUnicodeCharacter{2115}{\ensuremath{\mathbb{N}}}  % ℕ
\DeclareUnicodeCharacter{2102}{\ensuremath{\mathbb{C}}}  % ℂ
\DeclareUnicodeCharacter{211D}{\ensuremath{\mathbb{R}}}  % ℝ
\DeclareUnicodeCharacter{2124}{\ensuremath{\mathbb{Z}}}  % ℤ
\DeclareUnicodeCharacter{211A}{\ensuremath{\mathbb{Q}}}  % ℚ
\DeclareUnicodeCharacter{2192}{\ensuremath{\to}}         % →
\DeclareUnicodeCharacter{21D2}{\ensuremath{\Rightarrow}} % ⇒
\DeclareUnicodeCharacter{2200}{\ensuremath{\forall}}     % ∀
\DeclareUnicodeCharacter{2203}{\ensuremath{\exists}}     % ∃
\DeclareUnicodeCharacter{2208}{\ensuremath{\in}}         % ∈
\DeclareUnicodeCharacter{2209}{\ensuremath{\notin}}      % ∉
\DeclareUnicodeCharacter{2227}{\ensuremath{\wedge}}      % ∧
\DeclareUnicodeCharacter{2228}{\ensuremath{\vee}}        % ∨
\DeclareUnicodeCharacter{2229}{\ensuremath{\cap}}        % ∩
\DeclareUnicodeCharacter{222A}{\ensuremath{\cup}}        % ∪
\DeclareUnicodeCharacter{2245}{\ensuremath{\cong}}       % ≅
\DeclareUnicodeCharacter{2248}{\ensuremath{\approx}}     % ≈
\DeclareUnicodeCharacter{2260}{\ensuremath{\neq}}        % ≠
\DeclareUnicodeCharacter{2264}{\ensuremath{\leq}}        % ≤
\DeclareUnicodeCharacter{2265}{\ensuremath{\geq}}        % ≥
\DeclareUnicodeCharacter{22C5}{\ensuremath{\cdot}}       % ⋅
\DeclareUnicodeCharacter{2295}{\ensuremath{\oplus}}      % ⊕
\DeclareUnicodeCharacter{2297}{\ensuremath{\otimes}}     % ⊗

% ---- Mathematics & symbols ------------------------------------------
\usepackage{amsmath,amssymb,bm,mathrsfs,mathtools}
\usepackage{amsthm}             % theorem environments (load BEFORE hyperref)

% ---- Theorem environments (shared counter set, TECT canonical) -------
% All theorems / lemmas / propositions / corollaries / remarks share a
% single counter so cross-references read consistently across a paper.
% Definitions get a separate counter (definitions are numbered in their
% own sequence by editorial convention).
\newtheorem{theorem}{Theorem}
\newtheorem{lemma}[theorem]{Lemma}
\newtheorem{proposition}[theorem]{Proposition}
\newtheorem{corollary}[theorem]{Corollary}

\theoremstyle{remark}
\newtheorem{remark}[theorem]{Remark}
\newtheorem{example}[theorem]{Example}

\theoremstyle{definition}
\newtheorem{definition}[theorem]{Definition}
\newtheorem{conjecture}[theorem]{Conjecture}
\newtheorem{hypothesis}[theorem]{Hypothesis}

% ---- Tables / floats / graphics --------------------------------------
\usepackage{booktabs}           % \toprule / \midrule / \bottomrule
\usepackage{array}
\usepackage{graphicx}
\usepackage{enumitem}           % \begin{itemize}[leftmargin=...] options
\usepackage{hyphenat}           % \hyp{} for breakable hyphens in \texttt
\usepackage{seqsplit}           % \seqsplit{} for long unbreakable strings
\sloppy                         % allow stretchier inter-word spacing to
                                % reduce overfull \hbox warnings on long URLs
                                % and long \texttt tokens (paragraph-level).
\emergencystretch=3em           % last-resort hbox stretch budget

% ---- Cross-references (hyperref last among the standard packages,
%      cleveref AFTER hyperref) ----------------------------------------
\usepackage[%
  colorlinks=true,%
  linkcolor=blue,%
  citecolor=blue,%
  urlcolor=blue,%
  breaklinks=true,%
  bookmarks=true,%
  bookmarksnumbered=true%
]{hyperref}
\usepackage[capitalise,nameinlink]{cleveref}

% ---- TECT-canonical author block macros -----------------------------
% (defined here so every paper has the same author / affiliation style)
\newcommand{\tectauthor}{%
  \author{Jusang Lee}%
  \email{jtkor@outlook.com}%
  \affiliation{Independent researcher, Republic of Korea}%
}

% ---- TECT-canonical archive footer ----------------------------------
\newcommand{\tectarchivefooter}{%
  \par\medskip
  \noindent\small\textit{Canonical archive}:
  \url{https://github.com/TECT-OpenPhysics/TECT}.\par%
}

% ---- TECT TODO / AUDIT-FLAG / HONEST-SCOPE inline marks --------------
% Used in drafts to highlight pending audit items without disturbing
% the body text style. Suppress via \tectsuppressmarks (define empty).
\providecommand{\tectauditflag}[1]{\textcolor{red}{\textbf{[AUDIT-FLAG: #1]}}}
\providecommand{\tecttodo}[1]{\textcolor{orange}{\textbf{[TODO: #1]}}}
\providecommand{\tecthonestscope}[1]{\textit{Honest scope: #1.}}

% =====================================================================
% End of TECT canonical preamble
% =====================================================================


\begin{document}

\title{A provisional Kibble--Zurek cosmology programme note for TECT
\\ \normalsize{\textit{(NOT a canonical GAP-4 closure paper. The
canonical F-GAP4 shell per Math300 / Math311 is a defect-mass
verdict window $\mu_{\rm defect} \in [10^{13}, 10^{17}]\,
\mathrm{GeV}$, NOT an SKA/IPTA non-detection condition. The
predicted $\Omega_{\rm GW}$ peak frequency in the present draft is
internally inconsistent with PTA bands. External use FORBIDDEN at
the present canonical tier per Math314-AddD audit.)}}}

\author{Jusang Lee}
\email{jtkor@outlook.com}
\affiliation{Independent researcher, Republic of Korea}

\date{\today}

\begin{abstract}
This document is a PROVISIONAL Kibble--Zurek (KZ) cosmology programme
note for the Topological Energy Condensate Theory (TECT). The KZ
mechanism applied to the BCC topological condensate freeze-out
generates a spectrum of topological defects (candidate monopoles,
vortices, BCC domain walls) whose gravitational-wave signature is
sharply peaked at a frequency corresponding to the condensate shell
wavenumber. The original draft proposed a PROVISIONAL
$\Omega_{\rm GW}$ peak at $f \sim 10^{-16}$ Hz (incorrectly labelled
``pulsar-timing array band'') and re-defined the F-GAP4 gate as an
SKA/IPTA non-detection condition.

\textbf{Audit-flag (AUDIT-2026-05-02-Wave7, Math314-AddD)}: both
elements are inconsistent with current canonical bookkeeping and
with standard GW-observation phenomenology:
(a) the proposed $f \sim 10^{-16}$ Hz peak is below ALL existing GW
observational bands (PTA: nHz $\sim 10^{-9}$ Hz; LISA: mHz; LIGO:
kHz). The original ``PTA band'' label is incorrect; the paper later
cites SKA/IPTA sensitivity at $\sim 10^{-9}$ Hz, contradicting its
own peak-frequency prediction. The $\Omega_{\rm GW}$ peak frequency
must be re-derived from first principles before external citation.
(b) the canonical F-GAP4 shell per Math300 / Math311 is a
defect-mass verdict window
$\mu_{\rm defect} \in [10^{13},\,10^{17}]\,\mathrm{GeV}$ with
solver / grid certificates, NOT an SKA/IPTA non-detection condition.
Redefining F-GAP4 within this paper would change the canonical gate
without going through the canonical-source-hierarchy update process.
(c) ``Paper 16 completes the GAP-4 closure'' wording is incompatible
with the current GAP-4 = T3 / Stage-2 composite = T3 canonical tier.

External use of the present draft is FORBIDDEN at the present
canonical tier; the document is preserved as a candidate
KZ-cosmology programme note pending (i) re-derivation of the
$\Omega_{\rm GW}$ peak frequency, (ii) honest separation between this
candidate cosmological prediction and the canonical F-GAP4
defect-mass verdict shell, and (iii) demotion of the closure claim
to a candidate-prediction framing.
\end{abstract}

\maketitle

% =====================================================================
\section{Introduction --- provisional KZ-cosmology programme,
NOT a current canonical GAP-4 closure paper}\label{sec:intro}
% =====================================================================

Kibble--Zurek (KZ) formation of topological defects is a universal
mechanism in phase transitions: when a system is cooled through a
critical point faster than the microphysical relaxation timescale,
causally disconnected regions freeze out in different ground states
and topological defects appear at the boundaries between them. In
standard cosmology the mechanism is invoked in axion physics, in
beyond-SM monopole and cosmic-string scenarios, and in the
freeze-out of light-pseudo-scalar fluctuations.

TECT proposes a concrete realisation in the BCC topological
condensate: the condensate undergoes a phase transition as the
universe cools below the Brazovskii instability scale
$T_{\mathrm{crit}}\sim 10^{16}\,\mathrm{GeV}$, and the KZ mechanism
generates a spectrum of candidate defects (monopoles, vortices,
BCC domain walls) with characteristic wavenumber
$\sim q_0\sim \ell_P^{-1}$. These defects can in principle radiate
gravitational waves, and the resulting stochastic background
$\Omega_{\mathrm{GW}}(f)$ is the natural target of any
``cosmological observables'' programme for TECT.

\textbf{Honest scope of this paper.} The present document is a
\emph{provisional} KZ-cosmology programme note. It is \emph{not} a
current canonical GAP-4 closure paper. Two structural points must be
made up-front because earlier drafts mis-stated them:
\begin{itemize}
  \item The naive estimate of the peak frequency of the
    $\Omega_{\mathrm{GW}}$ spectrum at the TECT operating point sits
    far below all existing GW observational bands, including the
    pulsar-timing-array (PTA) band ($f \sim 10^{-9}\,\mathrm{Hz}$).
    The peak frequency derived in \S\ref{sec:fpeak} below is
    \emph{not} in the PTA band; the earlier ``directly observable
    via PTA'' wording is \emph{withdrawn} (Remark~\ref{rmk:not-PTA}).
  \item The canonical F-GAP4 falsification gate per Math300 /
    Math311 is a \emph{defect-mass verdict-window shell},
    $\mu_{\mathrm{defect}} \in [10^{13},\,10^{17}]\,\mathrm{GeV}$
    with PASS / FAIL / DEFER on solver and grid certificates.
    F-GAP4 is \emph{not} an SKA / IPTA non-detection condition; the
    SKA / IPTA gate of an earlier draft was a paper-internal
    redefinition and is \emph{withdrawn}
    (Remark~\ref{rmk:not-SKA-IPTA}).
\end{itemize}
Both withdrawals are recorded in this paper as honest scope
clarifications, and the F-GAP4 section below uses the canonical
defect-mass verdict-window shell, not the SKA/IPTA wording.

% =====================================================================
\section{The Kibble-Zurek defect spectrum}
% =====================================================================

\subsection{Cooling rate and freeze-out time}

The universe's expansion rate (Hubble parameter) during TECT-relevant epochs is:
\begin{equation}
H(T) = \frac{\dot{a}}{a} \sim \frac{T^2}{M_{\rm Planck}}
\end{equation}
for a radiation-dominated universe. The BCC condensate critical temperature is
(from the Brazovskii instability):
\begin{equation}
T_{\rm crit} \sim M_X,
\end{equation}
where $M_X$ is the characteristic shell energy (related to the wavenumber via
$M_X \sim q_0 c_T$).

The freeze-out temperature is determined by equating the relaxation timescale
$\tau_{\rm relax} \sim H^{-1}(T_{\rm freeze})$:
\begin{equation}
\tau_{\rm relax}(T_{\rm freeze}) \sim (T_{\rm crit} - T_{\rm freeze})^{1/2} \cdot |\dot{T}|^{-1},
\end{equation}
where $|\dot{T}| \sim H T$ is the cooling rate. Solving for $T_{\rm freeze}$ yields
(Math196):
\begin{equation}
T_{\rm freeze} \sim T_{\rm crit} \left(1 - \left(\frac{T_{\rm crit}}{M_{\rm Planck}}\right)^{1/4}\right).
\end{equation}

\subsection{Defect density and correlation length}\label{sec:KZ-defect-density}

The KZ prediction for defect density is:
\begin{equation}
n_{\rm defect} \sim \xi_{\rm freeze}^{-3}, \quad \xi_{\rm freeze} \sim (|\dot{T}| / T_{\rm crit}^2)^{-1/4}.
\end{equation}

Using the cosmological cooling rate $|\dot{T}| \sim H(T_{\rm freeze}) T_{\rm freeze}$:
\begin{equation}
\xi_{\rm freeze} \sim \left(\frac{M_{\rm Planck}}{T_{\rm freeze}^2}\right)^{1/4}.
\end{equation}

For the BCC case (Math159), the freeze-out correlation length at the TECT operating
point is:
\begin{equation}
\xi_{\rm freeze} \sim (10^3 \text{ to } 10^4) \, \ell_P,
\end{equation}
yielding a defect number density $n_{\rm defect} \sim (10^{-13} \text{ to } 10^{-15}) \, \ell_P^{-3}$
at that epoch.

\subsection{Naive peak-frequency estimate from defect dynamics}
\label{sec:fpeak}

Topological defects of characteristic size
$\xi_{\mathrm{freeze}}$ oscillate and radiate gravitational waves at
a frequency set by the inverse oscillation timescale,
\begin{equation}\label{eq:fGW}
  f_{\mathrm{GW}}^{(\mathrm{rest})}
  \;\sim\; \frac{1}{2\pi\,\tau_{\mathrm{osc}}},
  \qquad
  \tau_{\mathrm{osc}}\sim \xi_{\mathrm{freeze}} / c_T.
\end{equation}
In the observer frame the frequency is redshifted by the
inverse-temperature factor,
\begin{equation}\label{eq:redshift}
  f_{\mathrm{obs}}
  \;=\; f_{\mathrm{GW}}^{(\mathrm{rest})}\,
        \frac{a(T_{\mathrm{freeze}})}{a_0}
  \;=\; f_{\mathrm{GW}}^{(\mathrm{rest})}\,
        \frac{T_0}{T_{\mathrm{freeze}}},
\end{equation}
with $T_0 \approx 2.7\,\mathrm{K}$ the present CMB temperature.

For the BCC freeze-out at $T_{\mathrm{freeze}} \sim M_X
\sim 10^{16}\,\mathrm{GeV}$, the naive insertion of
$\xi_{\mathrm{freeze}}\sim (10^3\!-\!10^4)\,\ell_P$ from
\S\ref{sec:KZ-defect-density} gives, to within an order of magnitude,
\begin{equation}\label{eq:fobs-naive}
  f_{\mathrm{obs}}^{\mathrm{naive}}
  \;\sim\; 10^{-16}\;\text{to}\;10^{-15}\;\mathrm{Hz}.
\end{equation}

\begin{remark}[The estimate \eqref{eq:fobs-naive} is NOT in the PTA
band, and is below ALL existing GW observational bands]
\label{rmk:not-PTA}
The PTA band of the SKA / IPTA / NANOGrav programmes spans
$f\sim 10^{-9}\;\text{to}\;10^{-7}\,\mathrm{Hz}$. The naive estimate
\eqref{eq:fobs-naive} sits at $f \sim 10^{-16}\,\mathrm{Hz}$, that
is, \emph{seven orders of magnitude below the PTA lower edge}. It is
also far below the LISA band ($\sim$~mHz) and the LIGO band
($\sim$~kHz). The earlier-draft labelling of \eqref{eq:fobs-naive}
as ``in the PTA band'' is \emph{factually incorrect} and is
withdrawn here. The estimate \eqref{eq:fobs-naive} is closer to the
Hubble-scale band probed by CMB polarisation B-mode searches; even
there, the predicted amplitude $\Omega_{\mathrm{GW}}h^2\sim 10^{-15}$
of \S\ref{sec:OmegaGW} is well below current and projected
sensitivities. The estimate \eqref{eq:fobs-naive} is therefore
\emph{not currently observable in any GW band}; this is a programme
limitation, not a closure prediction.
\end{remark}

% =====================================================================
\section{Provisional $\Omega_{\mathrm{GW}}$-spectrum estimate
(programme-level only)}\label{sec:OmegaGW}
% =====================================================================

The gravitational-wave energy density per logarithmic frequency
interval is
\begin{equation}\label{eq:OmegaGW-def}
  \Omega_{\mathrm{GW}}(f)\,h^{2}
  \;=\;
  \int_{0}^{\infty} dt \int_{\text{modes}}
  \frac{d\rho_{\mathrm{GW}}}{df}(t,f).
\end{equation}
For KZ-formed defects the spectrum is peaked near the characteristic
frequency $f_{\mathrm{peak}}^{\mathrm{naive}}$ of \S\ref{sec:fpeak}
with a bandwidth $\Delta\log f \sim 1$ (defect-type dependent).
Math172 \cite{Math172} computes the detailed spectrum at the TECT
operating point
$(q_0, M_X, \lambda, \gamma, Y) = (14.32,\, 10^{16}\,\mathrm{GeV},\,
-0.0394,\, -0.0188,\, 280.2)$ and gives the \emph{provisional}
estimate
\begin{equation}\label{eq:OmegaGW-num}
  \Omega_{\mathrm{GW}}(f_{\mathrm{peak}})\,h^{2}
  \;\sim\; 10^{-15}
  \qquad
  \text{at } f \sim 10^{-16}\,\mathrm{Hz}.
\end{equation}

\begin{remark}[The estimate \eqref{eq:OmegaGW-num} is provisional and
NOT a closure prediction]\label{rmk:provisional-OmegaGW}
\eqref{eq:OmegaGW-num} is reported as a programme-level estimate
only. It is sensitive to (i) the operating point (particularly the
coupling $\lambda$, which controls the defect core size),
(ii) the four-defect-repair verdict pending in Math218-AddA, and
(iii) the freeze-out cooling-rate normalisation Math196 / Math159.
The estimate is \emph{not} part of any current canonical
falsification gate; per Remark~\ref{rmk:not-PTA}, the predicted
peak frequency is below all existing GW observational bands and
cannot be tested by any current or near-future
$\Omega_{\mathrm{GW}}$ measurement.
\end{remark}

% =====================================================================
\section{F-GAP4: the canonical defect-mass verdict-window shell
(NOT an SKA/IPTA non-detection condition)}\label{sec:FGAP4}
% =====================================================================

\begin{remark}[F-GAP4 is the Math300/Math311 defect-mass verdict
shell, NOT an SKA / IPTA non-detection condition]\label{rmk:not-SKA-IPTA}
The canonical F-GAP4 falsification gate per Math300 / Math311 is the
defect-mass verdict-window shell
\begin{equation}\label{eq:FGAP4}
  \boxed{\;
  \text{F-GAP4}: \quad
  \mu_{\mathrm{defect}}
  \;\in\;
  \big[\,10^{13}\,\mathrm{GeV},\;10^{17}\,\mathrm{GeV}\,\big]
  \;}
\end{equation}
extracted from the lowest non-trivial eigenvalue
$\lambda_{\min}(\Box_{E})$ of the candidate transverse Box operator
$\Box_{E}$ on the defect manifold, with PASS / FAIL / DEFER on
the canonical solver and grid certificates registered in Math300 and
the verdict-shell consumption tracked in Math311. F-GAP4 is
\emph{not} an SKA / IPTA non-detection condition; the
SKA / IPTA wording of an earlier draft was a paper-internal
redefinition that conflicted with the canonical-source hierarchy
(the internal canonical-source hierarchy). The earlier wording is \emph{withdrawn}.
\end{remark}

\paragraph{Status of the four current GW observational bands relative
to the present provisional estimate.}
The present estimate \eqref{eq:OmegaGW-num} is provided for
completeness only; \emph{none} of the current GW observational bands
overlaps with the predicted peak frequency:
\begin{itemize}
  \item PTA band (SKA / IPTA / NANOGrav): $f\sim 10^{-9}\,\mathrm{Hz}$
    --- \emph{seven orders of magnitude above} the predicted peak.
  \item LISA band: $f \sim 10^{-3}\,\mathrm{Hz}$ ---
    \emph{thirteen orders of magnitude above} the predicted peak.
  \item LIGO / Virgo band: $f \sim 10^{2}\,\mathrm{Hz}$ ---
    eighteen orders of magnitude above.
  \item CMB B-mode polarisation (Hubble-scale): the closest band in
    frequency, but the predicted amplitude
    $\Omega_{\mathrm{GW}}h^2 \sim 10^{-15}$ is well below current
    and projected B-mode sensitivities.
\end{itemize}
The earlier-draft assertion that the present prediction is
``directly observable via pulsar-timing arrays'' is therefore
\emph{factually wrong} and is withdrawn together with the F-GAP4
SKA / IPTA redefinition.

% =====================================================================
\section{Relation to Stage-2 and other pillars}\label{sec:stage2}
% =====================================================================

\textbf{Honest status.} The present provisional KZ-cosmology
programme note does \emph{not} complete a GAP-4 closure. The current
canonical GAP-4 tier per Math307 / Math310 is \textbf{T3 \textsc{proof
sketch}}, with promotion gated by the canonical defect-mass
verdict-window shell \eqref{eq:FGAP4} of Math300 / Math311 (PASS /
FAIL / DEFER on solver and grid certificates), \emph{not} by any
SKA / IPTA non-detection. The earlier-draft wording ``Paper~16
completes the GAP-4 closure'' is incompatible with the current
canonical Stage-2 composite tier $\min(T6,T6,T4,T3)=T3$ and is
\emph{withdrawn}.

\paragraph{What an F-GAP4 PASS would entail.}
A PASS on the canonical defect-mass verdict-window shell
\eqref{eq:FGAP4} would advance Pillar~4 sub-task~3 (the defect-mass
prediction) to T6 \textsc{proved conditional} and would simultaneously
support the BCC freeze-out cooling profile of Math159 / Math196 at
$T_{\mathrm{freeze}} \sim 10^{16}\,\mathrm{GeV}$. It would
\emph{not} by itself certify the provisional
$\Omega_{\mathrm{GW}}$ estimate of \eqref{eq:OmegaGW-num}, which is
unobservable in the near term per Remark~\ref{rmk:not-PTA}.

\paragraph{What an F-GAP4 FAIL would entail.}
A FAIL on \eqref{eq:FGAP4} would require
(i) revision of the KZ rescope of Math159 / Math196 (scope narrowing
or $T$-scaling modification);
(ii) re-evaluation of the defect-mass prediction in Pillar~4
sub-task 3; and (iii) possible reversion of Stage-2 composite-tier
claims dependent on the defect-mass shell.

% =====================================================================
\section{Conclusion --- provisional programme note, NOT a closure
paper}\label{sec:conclusion}
% =====================================================================

The KZ defect-formation mechanism in TECT generates a candidate
stochastic gravitational-wave background whose order-of-magnitude
peak-frequency estimate sits at $f \sim 10^{-16}\,\mathrm{Hz}$ ---
seven or more orders of magnitude below all current GW observational
bands (Remark~\ref{rmk:not-PTA}). The accompanying amplitude
estimate $\Omega_{\mathrm{GW}}h^2 \sim 10^{-15}$ is provided only as
a programme-level estimate and is \emph{not} part of any current
canonical falsification gate (Remark~\ref{rmk:provisional-OmegaGW}).
The canonical F-GAP4 gate is the defect-mass verdict-window shell
$\mu_{\mathrm{defect}} \in [10^{13},\,10^{17}]\,\mathrm{GeV}$ of
\eqref{eq:FGAP4}, with PASS / FAIL / DEFER on solver / grid
certificates per Math300 / Math311; it is \emph{not} an SKA / IPTA
non-detection condition (Remark~\ref{rmk:not-SKA-IPTA}).

The present document is therefore preserved as a \emph{provisional
KZ-cosmology programme note} pending (i) re-derivation of the
$\Omega_{\mathrm{GW}}$ peak frequency from a more careful defect-
dynamics calculation; (ii) honest separation between the cosmological
prediction and the canonical defect-mass verdict-window F-GAP4 gate;
and (iii) the four-defect-repair verdict tracked in Math218-AddA
(deadline 2026-06-15). External use is FORBIDDEN at the present
canonical tier per Math314-AddD audit.

\begin{thebibliography}{99}

\bibitem{TECTRepo}
J. Lee, \textit{Topological Energy Condensate Theory (TECT) Repository},
\texttt{https://github.com/...} (2026).

\bibitem{Math156}
J. Lee, ``Round V1–V5 comprehensive audit verdict,''
\textit{Math156}, 2026-04-25.

\bibitem{Math159}
J. Lee, ``KZ rescope from inflation,''
\textit{Math159}, 2026-04-25.

\bibitem{Math172}
J. Lee, ``PROVISIONAL prediction $\Omega_{\mathrm{GW}} \sim 10^{-15}$ at PTA-band frequency,''
\textit{Math172}, 2026-04-25.

\bibitem{Math196}
J. Lee, ``KZ rate Friedmann coupling,''
\textit{Math196}, 2026-05-01.

\bibitem{Math218AddA}
J. Lee, ``E3' cosmological realization four-defect repair,''
\textit{Math218-AddA}, 2026-04-29.

\end{thebibliography}

\end{document}
